\chapter*{Abstract}
\phantomsection\addcontentsline{toc}{chapter}{\protect Abstract}
\thispagestyle{plain}

Organic semiconductors are promising materials for lightweight, flexible, and
large-area electronic devices. Their performance, however, cannot be described
by material properties alone. In organic thin-film transistors, measured currents
are shaped by the semiconductor, the device geometry, the contacts, the
dielectric interface, and time-dependent charging processes. This makes mobility
extraction intrinsically model-dependent and motivates measurement strategies
that separate sheet transport in the semiconductor from parasitic device
effects.

This thesis investigates this problem using bottom-gate C8-BTBT organic
field-effect transistors and a current-driven gated van der Pauw approach. The
work connects fabrication, device-quality checks, conventional transistor
measurements, four-terminal sheet-conductance analysis, and transient circuit
tests. The central aim is not only to demonstrate field-effect operation, but to
evaluate how reliably transport parameters can be obtained when contacts,
leakage paths, traps, and bias history may influence the electrical response.

Conventional current--voltage measurements confirm that the fabricated devices
operate as $p$-type organic transistors with clear gate modulation and low gate
leakage. At the same time, their gradual turn-on, finite background resistance,
and nonideal transfer characteristics show that two-terminal measurements remain
strongly affected by contact and injection effects. The gated van der Pauw
measurements reduce this ambiguity by probing the accumulated film through a
four-terminal geometry. The resulting sheet conductance is reproducible for
different imposed currents and yields an effective thin-film mobility of about
\SI{3.1}{\centi\metre\squared\per\volt\per\second}.

Transient and amplifier measurements further show that the device response is
not purely static. Slow relaxation and circuit time constants affect the dynamic
behaviour and must be considered when interpreting transfer curves or designing
organic transistor circuits. Overall, the thesis demonstrates that reliable
characterization of organic field-effect devices requires a combination of
careful fabrication control, contact-independent transport analysis, and
awareness of time-dependent effects. The gated van der Pauw method is therefore
shown to be a valuable tool for extracting robust transport parameters from
organic thin-film transistors.
